Multi-hop questions often fail in retrieval-augmented generation (RAG) because the context contains the endpoints of a reasoning chain but not the bridge passage that links them. We call this \emph{bridge failure} and measure its cost: across four generators from two vendors, answer F1 is 72.6--79.9\% when the context contains all gold passages and 19.2--24.1\% when it does not, and the generators abstain on 53--63\% of bridge-failed contexts. We propose SES-RAG, a greedy context selector that makes no LLM call and adds to each candidate's relevance the change in the algebraic connectivity (Fiedler value) of the selected evidence subgraph under a token budget. On 861 HotpotQA bridge questions that were not used for tuning, it retrieves a complete context about 10 points more often than Hybrid RAG and improves F1 by 4.0--6.5 points for every generator; on 2WikiMultiHopQA it improves gold-complete retrieval by 8--19 points over all LLM-free baselines without re-tuning. The evidence also sets clear limits. SES-RAG is not significantly better than the strongest graph baselines on HotpotQA, ties them on MuSiQue, and trails an LLM-driven iterative retriever by 11 points of gold-complete retrieval and 7--10 points of F1 on HotpotQA (22 points of gold-complete retrieval on 2WikiMultiHopQA). Simpler connectivity rewards (component count, unnormalised $\lambda_2$) are statistically indistinguishable from the spectral one after multiple-comparison correction, its advantage over the same selector without the connectivity term disappears when the graph lacks title-mention edges, and combining iterative retrieval with the same packing rule gives no significant gain on HotpotQA and a loss on MuSiQue. We release code, frozen parameters and all model outputs.
